% ======================================================================
\section{Conclusion}\label{sec:conclusion}

This paper asked whether countries have met the basic social needs of their populations while staying within a fair per-person share of the climate boundary, and how this changed between \nFirstYear{} and \nLastYear. It combined four social thresholds with consumption-based \COtwo{} emissions for the \nNCons{} countries with such data, from a panel of \nNCountries{} countries, and compared each country's emissions with an equal-per-capita boundary of \nBoundary~t per person per year.

The answer has four parts. First, none of the \nNCons{} countries with consumption-based emissions data met the social foundation within the climate boundary in any year of the period, and the \nProdSafeYearsAll{} country-years in which production-based emissions place a country there disappear once imports are counted. Every one of the \nNCrossers{} countries that crossed the social foundation did so while overshooting the boundary. Second, the carbon requirement, the \COtwo{} per person at which half of countries meet the foundation, fell by about \nReqTrend\% a year, from roughly \nReqStart~t in the early 1990s to \nReqEnd~t in \nLastYear{}. Meeting basic needs has become much less carbon-intensive. Yet the requirement remained \nReqEndMultiple{} times the boundary, and its exact level depends on how the social foundation is defined. Third, absolute decoupling of emissions from income became common in the 2010s, but almost no country is decoupling fast enough to reach the boundary by mid-century. Even the fastest rate any decoupling country achieved would bring only a minority of them there in time. Fourth, the configurations of shortfall and overshoot are persistent, and the historical path into the social foundation has run through overshoot.

Together these findings point to a clear conclusion. The relationship between emissions and basic social needs is not fixed, and it has improved substantially. A few countries have shown that the four social thresholds can be met with emissions close to the boundary. But no country in the sample has yet combined meeting them with emissions within the boundary on a consumption basis. The current pace of change in both the carbon cost of development and the decarbonisation of countries above the boundary falls far short of what the climate boundary requires. Closing the gap will require emission cuts much faster than any country has so far sustained, together with ways of providing for basic needs that keep their carbon cost falling. These conclusions concern one planetary boundary, climate change, and four social indicators. They do not cover the other planetary boundaries or the many dimensions of well-being the four indicators leave out.

% ======================================================================
\section*{Data and code availability}
All data used in this study are publicly available. The Global Carbon Budget, Maddison Project and greenhouse-gas data were obtained from the public repository of \citet{owid2025co2} on 20~October~2026. The social indicators come from the World Bank's World Development Indicators. The code, the raw input files with SHA-256 checksums, the processed panel and all results will be released in a public repository on publication. Running a single command reproduces every number, table and figure in this paper.

\section*{Acknowledgements}
I thank Dr.\ Kamal Kumar Murari and Prof.\ Unmesh Patnaik for their guidance during my undergraduate studies at the Tata Institute of Social Sciences, from which this research developed. Any errors are my own.

\section*{Funding}
This research received no specific funding.

\section*{Declaration on the use of AI tools}
% AUTHOR: review and adapt this statement to your institution's and the journal's policy before submission.
An AI assistant (Claude, Anthropic) was used to help write analysis code, run the analyses and draft text. The author designed the study, reviewed all code, results and text, and takes full responsibility for the content.

\section*{Declaration of competing interests}
The author declares no competing interests.
